\chapter{三维热模型与 32 阶热降阶}
\label{ch:heat}

\section{热量从哪里产生，又往哪里去}
第~\ref{ch:em} 章得到海水的空间焦耳热 $q_{\rm sea}(\mathbf x,t)$ 和两只铜线圈的损耗功率 $P_{{\rm cu},1},P_{{\rm cu},2}$。这些是热方程的输入。铜发热先使导线升温，再经过绝缘层进入封装；海水中的感应电流直接加热海水；封装与海水之间还会交换热量。即使断电，温度差依然推动热量传递。

\begin{figure}[htbp]
\centering\resizebox{0.95\textwidth}{!}{\begin{tikzpicture}[node distance=12mm and 12mm]
\node[therm,text width=4.4cm] (sea) {海水三维 FV\\含封装附近的八层热层};
\node[therm,above left=of sea,text width=3.3cm] (p1) {发射端封装\\三维温度场};
\node[therm,above right=of sea,text width=3.3cm] (p2) {接收端封装\\三维温度场};
\node[drive,above=of p1,text width=3.3cm] (c1) {发射线圈铜\\铜损耗 $P_{\rm cu,1}$};
\node[drive,above=of p2,text width=3.3cm] (c2) {接收线圈铜\\铜损耗 $P_{\rm cu,2}$};
\draw[arr,<->] (c1)--node[right,font=\scriptsize] {绝缘传热} (p1);
\draw[arr,<->] (c2)--node[left,font=\scriptsize] {绝缘传热} (p2);
\draw[arr,<->] (p1)--(sea);
\draw[arr,<->] (p2)--(sea);
\node[font=\footnotesize,align=center,text=BFink,below=5mm of sea] {海水由 $q_{\rm sea}(\mathbf x,t)$ 直接加热，封装由铜和海水传热\\降阶后保留这五个分区的储热收支};
\end{tikzpicture}
}
\caption{五个热学分区及其主要传热联系。实际空间求解由许多有限体积节点组成。}
\label{fig:heat}
\end{figure}

这里称为“五材料”热 ROM，具体是五个守恒分区：海水、发射封装、接收封装、发射铜和接收铜。铜与封装分别属于相同材料类别，但两侧的热量和温度需要单独记账。

\section{热容量、导热和散热的直观理解}
给一块物体输入相同功率，热容量较大时升温较慢；与周围物体的导热连接越强，热量越容易流向外部。这对应热方程中的三组物理量：
\begin{itemize}
\item \textbf{热容量}：储存一定热量时温度能够升高多少；
\item \textbf{导热系数与界面热阻}：不同位置之间交换热量的快慢；
\item \textbf{边界换热}：热量如何传到更远的海水环境。
\end{itemize}
对一个材料区域，连续方程为
\[
\rho c\,\frac{\partial T}{\partial t}
=\nabla\cdot(k\nabla T)+q.
\]
左边是储热变化，第一项是空间热扩散，$q$ 是内部热源。两材料接触处通过接触传热或连续热通量相连。

\section{为什么海水要有封装附近的细热层}
原 BFZI 的海水有限体积单元可能达到厘米量级。封装外表面的温度梯度集中在更小的距离上，用整块海水的平均温度近似封装附近的温度，会使界面换热依赖于网格大小。

因此物理参考增加了贴附封装表面的\textbf{有限深度海水热层}，按当前方案采用约 2\,mm 厚、八层离散；海水热层与外部海水体积连接。每层具有自己的热容量和导热连接，避免把近表面的热滞后压缩为一个任意固定传热系数。

在已有 P03/P07 几何中，典型完整热状态由约 1312 个海水湿 FV、64 个封装 FV、两个铜节点以及 1024 个近表面海水层节点组成，总计 2402 个热状态。新的验证样本使用其他原生网格数量，状态维数可以随之改变。

当前铜线圈温度采用\textbf{两只线圈各一个等效铜热节点}；海水与双封装具有三维空间分布，近表面层还记录更细的海水热扩散。

\section{热参数如何进入真实方程}
用户可以逐工况给定绝缘导热系数、绝缘厚度、封装导热系数以及封装与海水之间的换热系数。例如，铜与封装之间的导热连接由绝缘层厚度、材料导热系数和实际绕组接触几何共同决定。

这样的参数改变会进入热算子，而不是只作为神经网络中的一个标签。改变绝缘厚度会改变铜向封装释放热量的时间过程；改变封装导热系数会改变封装内部温度分布；改变海水换热能力会改变封装与外部环境之间的热交换。

\section{将三维热方程写成矩阵}
有限体积离散后的温升向量记为
$\boldsymbol\theta=T-T_{\rm amb}$，有
\begin{equation}
\label{eq:fullheat}
\mathbf C_T\dot{\boldsymbol\theta}
+\mathbf H_T\boldsymbol\theta
=\mathbf Q(t).
\end{equation}
$\mathbf C_T$ 是热容量对角矩阵，$\mathbf H_T$ 包括体内导热、材料间热交换和边界散热，$\mathbf Q(t)$ 是当前空间热功率。

一次隐式欧拉热步可写成
\begin{equation}
\label{eq:implicit}
(\mathbf C_T+\Delta t\,\mathbf H_T)
\boldsymbol\theta_{n+1}
=\mathbf C_T\boldsymbol\theta_n+\Delta t\,\mathbf Q_{n+1}.
\end{equation}
之所以使用隐式形式，是因为铜、封装和海水具有相差很大的热时间尺度，显式推进可能需要非常小的时间步。每一步的热源又依赖当前阻抗与电流，因此代码会进行有限次热源—温度迭代。

\section{什么是热降阶}
假设一个具体几何有 2402 个温度未知数，若每一步都求解 2402 维方程，单条长历史需要重复很多次矩阵运算。热降阶的思路是：系统真正经历的温度分布往往由少量主要形状组合而成。

例如，一个模式代表整体海水升温，一个模式代表接收封装局部热点，另一个模式代表封装到海水的温度梯度。若用 32 个模式就能近似大多数重要温度分布，那么只需保存 32 个随时间变化的系数。

写成数学式就是
\[
\boldsymbol\theta(t)\approx \Phi\mathbf a(t),\qquad
\Phi\in\R^{N_T\times 32},\quad \mathbf a(t)\in\R^{32}.
\]
$\Phi$ 的每一列是一个固定的空间温度形状，$\mathbf a(t)$ 决定在某个时刻该形状占多大比重。$\Phi$ 是空间坐标的表示，$\mathbf a$ 是动态状态。

\section{32 个模式如何生成}
当前代码 \code{geom\_rom.py} 的构基步骤是：
\begin{enumerate}
\item 从冷态物理锚读取真实热容量 $\mathbf C_T$、导热算子 $\mathbf H_T$ 和五类电磁热源；
\item 对海水双端口损耗的 $11,12,22$ 组合以及两只铜的单位加热，求多个短时、长时和冷却热响应；
\item 将响应按热容量加权，提取具有代表性的温度形状；
\item 强制保留五个热学分区的常数温度模式；
\item 对其余空间变化进行奇异值分解，保留主要模态，形成总计最多 32 个模式。
\end{enumerate}

这套热基来自\textbf{真实冷态传热算子与物理热源响应}，构基时没有使用目标留出工况的动态温度真值。热容量加权正交关系为
\[
\Phi^\mathsf T\mathbf C_T\Phi=I.
\]
代入式~\eqref{eq:fullheat} 并投影，得到
\begin{equation}
\label{eq:romheat}
\dot{\mathbf a}
+\mathbf K_r\mathbf a
=\Phi^\mathsf T\mathbf Q,
\qquad
\mathbf K_r=\Phi^\mathsf T\mathbf H_T\Phi.
\end{equation}
在线只需求解小型矩阵方程，然后通过 $\Phi\mathbf a$ 将温度恢复至原三维 FV 单元。

\section{降阶后怎样保留能量收支}
为什么一定要保留五个材料常数模式？因为它们分别代表五个分区的积分热量。保留它们，就能将每个分区的储热与传入传出热流在 ROM 中正确对应。

整个系统的总热收支可写为
\[
\frac{\dd E_{\rm stored}}{\dd t}
=P_{\rm sea}+P_{{\rm cu},1}+P_{{\rm cu},2}
-P_{\rm out},
\]
其中 $P_{\rm out}$ 是向环境传出的净热功率。对铜、封装、海水等单独分区，还要计入分区之间的界面热流。界面流出一个分区的热，必须以相同数值进入另一个分区。

32 阶投影之后，逐个 FV 小单元的传热方程可能含有降阶截断误差；五个分区的总能量预算仍可直接计算并检查。仓库的 \code{GeomHeatROM} 会记录每个分区的能量收支。

\section{为什么既有热 ROM，又要神经网络}
两者解决的是不同计算问题。

热 ROM 用来回答：\emph{已知热源，热量下一步怎样流动？} 这件事受热容量和导热方程严格约束，直接用 ROM 是清晰、稳定的方法。

神经网络用来回答：\emph{经历此前的运行历史后，海水电磁反作用在新时刻如何变化？} 这个量如果每步都重新用 BFZI 电磁求解器计算，成本较高，因此用因果记忆和物理切线共同加速。

热 ROM 提供温度反馈，神经网络提供动态电磁先验，物理切线提供高精度电磁温度校正。完整连接在第~\ref{ch:coupling} 章给出。

\section{代码与历史结果}
主要代码为 \code{sdfmpneo\_bfzi/geom\_rom.py} 和 \code{heat\_rom.py}。原 B 阶段对 P03/P07 建立了 2402$\to$32 阶 ROM，并在独立激励上通过温度与动态阻抗验收。那一阶段的总在线速度仍受逐热步 BFZI 电磁求解影响。后续加入快速电磁代理之后，完整时程才得到显著加速。
